\chapter{传统机器学习模型}

机器学习模型的种类繁多，但从学习方式来看，可以先沿两条主线区分。一条是\textbf{监督学习（Supervised Learning）}，训练数据包含特征和对应标签，模型目标是学习从特征到标签的映射函数；另一条是\textbf{无监督学习（Unsupervised Learning）}，训练数据只有特征而没有标签，模型目标是发现数据内部的结构。

监督学习又可以根据标签类型进一步细分。当标签是连续数值时，任务称为\textbf{回归（Regression）}，例如预测房价、预测交通流量；当标签是有限个离散类别时，任务称为\textbf{分类（Classification）}，例如判断邮件是否为垃圾邮件、识别图像中的物体。回归与分类共享同一类学习范式，都需要带标签的数据、都需要定义损失函数、都通过优化参数来提升预测能力，区别仅在于输出的类型与对应的损失函数形式。

本章关注的重点不在于手工实现每一种算法，而在于理解每个模型的核心假设、适用条件与常见失效模式，从而在 AI 协作建模时能够判断推荐方案是否合适、验证生成结果是否可靠。

\section{回归模型}

回归问题旨在预测连续数值型输出。训练数据形如 $\{(\mathbf{x}_i, y_i)\}_{i=1}^N$，其中 $y_i \in \mathbb{R}$，模型目标是学习一个函数 $f(\mathbf{x};\theta)$，使得在新样本上的预测 $\hat{y} = f(\mathbf{x};\theta)$ 尽量接近真实值。

在回归任务中，模型好坏通常不能只看一个指标。常见指标包括 MAE、MSE/RMSE、MAPE 和 $R^2$。其中，MAE 直接衡量平均绝对误差，单位与原始标签相同；RMSE 会放大较大的预测偏差，适合关注大误差代价较高的场景；MAPE 衡量相对误差，定义为
\[
\mathrm{MAPE}=\frac{1}{N}\sum_{i=1}^{N}\left|\frac{y_i-\hat{y}_i}{y_i}\right|\times100\%.
\]
例如，MAPE 为 $8\%$ 可以理解为预测值平均偏离真实值约 $8\%$。它在需求预测、销量预测和交通流量预测中很常用，因为这些任务往往关心"相对偏差有多大"。不过，当真实值 $y_i$ 接近 0 或等于 0 时，MAPE 会变得不稳定，因此使用时需要结合数据范围判断是否合适。

\subsection{线性回归}

线性回归是最简单也是最基础的回归模型，假设输入特征与输出变量之间存在线性关系。对于单个特征 $X$ 和输出 $Y$，模型为
\[
Y = \beta_0 + \beta_1 X + \epsilon,
\]
其中 $\beta_0$ 是截距，$\beta_1$ 是斜率，$\epsilon$ 是误差项。推广到多个特征 $X_1, X_2, \ldots, X_n$ 时，模型变为多元线性回归：
\[
Y = \beta_0 + \beta_1 X_1 + \beta_2 X_2 + \ldots + \beta_n X_n + \epsilon.
\]
模型的目标是找到最优的 $\beta_0, \beta_1, \ldots, \beta_n$，通常通过最小化残差平方和（Residual Sum of Squares, RSS）来实现：
\[
\min_{\beta_0, \ldots, \beta_n} \sum_{i=1}^{m} \left(y_i - (\beta_0 + \beta_1 x_{i1} + \ldots + \beta_n x_{in})\right)^2.
\]
线性回归简单、易于理解和解释，是许多更复杂模型的基石。然而，它要求特征与目标之间存在线性关系；如果这一假设不成立，模型性能会受到明显限制。

\paragraph{实现占位}
我们将使用 \texttt{scikit-learn} 自带的 diabetes 数据集演示 \texttt{LinearRegression}。该数据集来自真实医学观测，目标变量是一年后的疾病进展指标；这里选取体重指数（BMI）作为单一解释变量，便于直观看到线性回归如何拟合连续数值。

\begin{quote}
\textbf{【实现占位】} 本段原有代码不直接迁入；后续案例将以 AI 协作过程、关键决策和验证结果为主。
\end{quote}

\subsection{多项式回归}

当特征与目标变量之间存在非线性关系时，简单的线性回归可能无法很好地捕捉这种模式。多项式回归通过在原始特征中添加多项式特征（如 $X^2, X^3$ 等）来扩展线性模型，从而能够拟合非线性的数据。本质上，多项式回归仍然是线性回归的一种形式，因为它对变换后的特征是线性的。

对于单个特征 $X$ 和 $Y$，一个 $d$ 次多项式回归模型表示为：
\[
Y = \beta_0 + \beta_1 X + \beta_2 X^2 + \ldots + \beta_d X^d + \epsilon.
\]
选择合适的阶数非常重要：阶数过高可能导致过拟合，阶数过低可能导致欠拟合。

\paragraph{实现占位}
我们继续使用 diabetes 数据集中的 BMI 特征，加入二次项，观察多项式特征是否能捕捉 BMI 与疾病进展之间的非线性关系。

\begin{quote}
\textbf{【实现占位】} 本段原有代码不直接迁入；后续案例将以 AI 协作过程、关键决策和验证结果为主。
\end{quote}

\subsection{正则化线性模型：岭回归、Lasso 与弹性网}

普通线性回归在特征数量较多、特征之间存在多重共线性或训练样本相对较少时，容易出现参数估计不稳定或过拟合的问题。正则化线性模型在损失函数中添加惩罚项，通过约束参数的大小来缓解这些问题。根据惩罚项形式的不同，正则化线性模型主要分为三种：岭回归、Lasso 和弹性网，分别对应 $L_2$ 惩罚、$L_1$ 惩罚及两者的结合。

\textbf{岭回归（Ridge Regression）}在普通最小二乘的损失函数中添加了 $L_2$ 惩罚项：
\[
J_{\text{Ridge}}(\beta) = \sum_{i=1}^{m} (y_i - \hat{y}_i)^2 + \alpha \sum_{j=1}^{n} \beta_j^2,
\]
其中 $\alpha \geq 0$ 是正则化参数。当 $\alpha$ 增大时，对系数的惩罚越大，模型系数会趋向于零，从而降低模型复杂度、减少过拟合风险。岭回归的 $L_2$ 惩罚项会缩小所有系数，但不会使任何系数完全变为零（除非 $\alpha$ 趋于无穷大），这意味着它不会进行特征选择。

从几何角度看，普通最小二乘的损失函数可以画成一组椭圆形等高线，最优解位于等高线的中心；而岭回归相当于要求参数向量落在一个 $L_2$ 球内。最终解出现在"损失等高线"和"参数约束球"相切的位置。由于 $L_2$ 球是光滑的，相切点通常不会正好落在坐标轴上，因此岭回归倾向于把所有系数一起缩小，而不是直接把某些系数压成 0。

% 【图示占位】源稿图示待按本书风格重绘。

相关图示 中的椭圆表示损失函数等高线，圆形或菱形表示参数约束范围。Ridge 的相切点多落在坐标轴之外，因此常表现为整体收缩；Lasso 的尖角更容易与等高线相切在坐标轴上，因此可能把某些特征系数直接压为 0。

\textbf{Lasso 回归（Least Absolute Shrinkage and Selection Operator）}采用 $L_1$ 惩罚：
\[
J_{\text{Lasso}}(\beta) = \sum_{i=1}^{m} (y_i - \hat{y}_i)^2 + \alpha \sum_{j=1}^{n} |\beta_j|.
\]
$L_1$ 惩罚具有"稀疏性"的特点：它能够将某些不重要特征的系数直接压缩为零，从而实现自动特征选择。从等高线图上看，Lasso 的 $L_1$ 约束区域像一个带尖角的菱形；损失函数等高线向外扩张时，很容易先碰到菱形的尖角，而这些尖角正好位于某些坐标轴上，意味着某些参数会被压缩为 0。这使得 Lasso 特别适合高维场景下的特征筛选。但当存在高度相关的特征时，Lasso 倾向于随机选择其中一个特征而忽略其他特征。

\textbf{弹性网（ElasticNet）}结合了岭回归和 Lasso 的优点，在损失函数中同时包含 $L_1$ 和 $L_2$ 惩罚项：
\[
J_{\text{ElasticNet}}(\beta) = \sum_{i=1}^{m} (y_i - \hat{y}_i)^2 + \alpha \rho \sum_{j=1}^{n} |\beta_j| + \frac{\alpha (1-\rho)}{2} \sum_{j=1}^{n} \beta_j^2,
\]
其中 $\alpha \geq 0$ 是总的正则化强度，$\rho \in [0, 1]$ 是 $L_1$ 与 $L_2$ 惩罚的混合比例。当 $\rho = 1$ 时，弹性网退化为 Lasso；当 $\rho = 0$ 时，退化为岭回归。弹性网在处理高维数据和多重共线性时兼顾了特征选择能力和估计稳定性，实践中通常通过交叉验证选择最优的 $\alpha$ 和 $\rho$ 值。

\paragraph{实现占位}
\texttt{scikit-learn} 中分别有 \texttt{Ridge}、\texttt{Lasso} 和 \texttt{ElasticNet} 模型可以直接使用。

\begin{quote}
\textbf{【实现占位】} 本段原有代码不直接迁入；后续案例将以 AI 协作过程、关键决策和验证结果为主。
\end{quote}

\subsection{逐步回归}

逐步回归是一种特征选择方法，通过迭代地添加或删除特征来构建回归模型，目标是找到一个最优的特征子集，在提供良好预测能力的同时保持模型的简洁性。主要有三种变体：\textbf{向前选择（Forward Selection）}从空集开始逐步添加最有用的特征；\textbf{向后消除（Backward Elimination）}从包含所有特征的模型开始逐步删除影响最小的特征；\textbf{逐步法（Stepwise）}则结合两者，每步都考虑添加或删除特征。决定特征是否加入或删除的标准通常基于统计显著性（如 $p$ 值）或信息准则（如 AIC、BIC）。

逐步回归是一种有效的特征选择策略，可以帮助构建更简洁、更具解释性的模型。然而，它的贪婪性质可能无法找到全局最优特征子集，并且在特征数量非常大时计算成本较高。在现代机器学习中，通常有更稳健的特征选择方法（如 Lasso、递归特征消除等）可以替代。

\paragraph{实现占位}
\texttt{scikit-learn} 没有内置逐步回归函数，但可以通过循环和评估指标手动实现向前选择过程。

\begin{quote}
\textbf{【实现占位】} 本段原有代码不直接迁入；后续案例将以 AI 协作过程、关键决策和验证结果为主。
\end{quote}

\section{分类模型}

分类问题旨在预测离散类别标签。训练数据形如 $\{(\mathbf{x}_i, y_i)\}_{i=1}^N$，其中 $y_i \in \{1, 2, \ldots, K\}$，模型目标是学习从特征到类别的映射函数。与回归不同，分类模型通常先输出属于各类别的概率，再根据阈值或最大概率原则做出最终类别判断。

\subsection{逻辑回归}

尽管名字中带有"回归"，逻辑回归实际上是一个广泛用于二分类的线性模型。它假设在给定观测值 $\mathbf{x}_n$ 的情况下，属于类别 $C_1$（$y_n=1$）的概率可以写作 Logistic 函数（即 Sigmoid 函数）：
\begin{equation}
    P(y_n=1|\mathbf{x}_n) = \sigma(w_0 + \mathbf{w}^{\textrm{T}} \mathbf{x}_n) = \frac{e^{w_0 + \mathbf{w}^{\textrm{T}} \mathbf{x}_n}}{1 + e^{w_0 + \mathbf{w}^{\textrm{T}} \mathbf{x}_n}}.
\end{equation}
这等价于 Logit 形式（对数几率）：
\begin{equation}
    \log \left( \frac{P(y_n=1|\mathbf{x}_n)}{1 - P(y_n=1|\mathbf{x}_n)} \right) = w_0 + \mathbf{w}^{\textrm{T}} \mathbf{x}_n.
\end{equation}
决策规则直观：计算概率 $P(y_n=1|\mathbf{x}_n)$，若大于 0.5 则预测类别 $C_1$；否则预测类别 $C_2$（$y_n=0$）。相关图示 展示了 Logistic 函数最重要的性质：无论线性输入取多大或多小，输出都会被压缩到 $0$ 和 $1$ 之间，因此可以被解释为概率。

% 【图示占位】源稿图示待按本书风格重绘。

\paragraph{应用实例：学习时间与考试通过率}

为了直观理解，考虑一个具体的例子。假设有一组 20 名学生，他们的复习时间（0 到 6 小时）与是否通过考试（0 或 1）的数据如下表所示。

\begin{table}[H]
    \centering
    \caption{学生复习时间与考试结果原始数据}
    \label{tab:exam_data}
    \resizebox{\textwidth}{!}{
    \begin{tabular}{@{}lcccccccccc@{}}
        \toprule
        \textbf{Hours} & 0.5 & 0.75 & 1 & 1.25 & 1.5 & 1.75 & 1.75 & 2 & 2.25 & 2.5 \\
        \textbf{Pass} & 0 & 0 & 0 & 0 & 0 & 0 & 1 & 0 & 1 & 0 \\
        \midrule
        \textbf{Hours} & 2.75 & 3 & 3.25 & 3.5 & 4 & 4.25 & 4.5 & 4.75 & 5 & 5.5 \\
        \textbf{Pass} & 1 & 0 & 1 & 0 & 1 & 1 & 1 & 1 & 1 & 1 \\
        \bottomrule
    \end{tabular}
    }
\end{table}

通过逻辑回归拟合，可以得到如下模型系数表：

\begin{table}[H]
    \centering
    \caption{逻辑回归模型拟合结果}
    \label{tab:logistic_results}
    \begin{tabular}{lcccc}
        \toprule
        & \textbf{Coefficient} & \textbf{Std. Error} & \textbf{$z$-value} & \textbf{$p$-value} \\
        \midrule
        Intercept ($w_0$) & -4.0777 & 1.7610 & -2.316 & 0.0206 \\
        Hours ($w_1$) & 1.5046 & 0.6287 & 2.393 & 0.0167 \\
        \bottomrule
    \end{tabular}
\end{table}

结果显示"学习时间"与通过概率显著相关（$p=0.0167 < 0.05$）。基于拟合出的参数 $w_0 = -4.0777$ 和 $w_1 = 1.5046$，可以预测不同学习时间下的通过概率：
\[
P(\text{Pass}|x) = \frac{1}{1 + e^{-(-4.0777 + 1.5046 \cdot x)}}.
\]

部分预测结果如下表所示：

\begin{table}[H]
    \centering
    \caption{不同学习时间下的预测通过率}
    \label{tab:prob_prediction}
    \begin{tabular}{cc}
        \toprule
        \textbf{Hours of study} & \textbf{Probability of passing exam} \\
        \midrule
        1 & 0.07 \\
        2 & 0.26 \\
        3 & 0.61 \\
        4 & 0.87 \\
        5 & 0.97 \\
        \bottomrule
    \end{tabular}
\end{table}

可以看到，学习时间为 3 小时时通过概率约为 0.61；增加到 5 小时后，通过概率高达 0.97。

\paragraph{参数估计与梯度下降}

逻辑回归的参数求解使用极大似然估计（MLE）。令 $z_n = P(y_n=1|\mathbf{x}_n)$，对于 $N$ 个样本，似然函数定义为所有样本被正确分类的联合概率：
\begin{equation}
    L(\mathbf{w}, w_0|\mathcal{X}) = \prod_{n=1}^{N} (z_n)^{y_n} (1 - z_n)^{1 - y_n}.
\end{equation}
为了便于计算，通常取负对数转化为最小化交叉熵误差：
\begin{equation}
    E(\mathbf{w}, w_0|\mathcal{X}) = - \sum_{n=1}^{N} \left[ y_n \log z_n + (1 - y_n) \log (1 - z_n) \right].
\end{equation}
由于上述误差函数没有闭式解，需要使用梯度下降法迭代求解。利用 Logistic 函数的导数 $\frac{\mathrm{d} z}{\mathrm{d} a} = z(1-z)$，可以推导出参数 $w_j$ 的梯度：
\begin{equation}
    \frac{\partial E}{\partial w_j} = \sum_{n=1}^N (z_n - y_n) x_{nj},
\end{equation}
相应的梯度下降更新规则为（其中 $\eta$ 为学习率）：
\begin{equation}
    w_j \leftarrow w_j - \eta \sum_{n=1}^N (z_n - y_n) x_{nj}.
\end{equation}
这一公式很优雅：权重的调整量与预测误差 $(y_n - z_n)$ 和输入特征 $x_{nj}$ 成正比。

\begin{algorithm}[H]
    \caption{逻辑回归梯度下降算法}
    \begin{algorithmic}[1]
        \State 初始化 $w_j \leftarrow \mathrm{rand}(-0.01,0.01)$，$j=0,\ldots,D$
        \While{未满足收敛条件}
            \For{$j=0,\ldots,D$}
                \State $\Delta w_j \leftarrow 0$
            \EndFor
            \For{$n=1,\ldots,N$}
                \State $z_n \leftarrow \sigma(w_0+\mathbf{w}^{T}\mathbf{x}_n)$
                \For{$j=0,\ldots,D$}
                    \State $\Delta w_j \leftarrow \Delta w_j+\eta (y_n-z_n)x_{nj}$
                \EndFor
            \EndFor
            \For{$j=0,\ldots,D$}
                \State $w_j \leftarrow w_j+\Delta w_j$
            \EndFor
        \EndWhile
    \end{algorithmic}
\end{algorithm}

\paragraph{实现占位}
我们将使用 \texttt{scikit-learn} 库中的 \texttt{LogisticRegression} 模型。

\begin{quote}
\textbf{【实现占位】} 本段原有代码不直接迁入；后续案例将以 AI 协作过程、关键决策和验证结果为主。
\end{quote}

\paragraph{多项逻辑回归与 Softmax 回归}

当分类问题的类别数量 $K>2$ 时，二分类逻辑回归需要推广为多项逻辑回归（Multinomial Logit Model），也常称为 Softmax 回归。它为每个类别 $k$ 定义一个线性得分：
\[
s_{nk}=w_{k0}+\mathbf{w}_{k}^{T}\mathbf{x}_n,\qquad k=1,\ldots,K.
\]
这些得分本身还不是概率，可以为正或为负，不要求相加为 1。Softmax 函数将它们转换为一组合法概率：
\begin{equation}
    z_{nk}=P(C_k | \mathbf{x}_n) = \frac{\exp(w_{k0} + \mathbf{w}_k^{\textrm{T}} \mathbf{x}_n)}{\sum_{j=1}^K \exp(w_{j0} + \mathbf{w}_j^{\textrm{T}} \mathbf{x}_n)}.
\end{equation}
Softmax 有三个重要性质：所有 $z_{nk}$ 都大于 0；所有类别概率之和为 1；原始得分越大的类别，转换后的概率通常也越大。预测时，模型选择概率最大的类别：$\hat{y}_n=\arg\max_{k} z_{nk}$。

对于标签采用独热编码的 $K$ 分类问题，多分类交叉熵损失函数为：
\begin{equation}
    E(\{\mathbf{w}_k, w_{k0}\}_k | \mathcal{X}) = - \sum_{n=1}^N \sum_{k=1}^K y_{nk} \log z_{nk}.
\end{equation}
由于 $y_{nk}$ 是独热向量，上式对每个样本实际上只惩罚真实类别对应的概率：$-\log z_{n,y_n}$。参数更新规则为：
\begin{equation}
    \Delta \mathbf{w}_j = \eta \sum_{n=1}^N (y_{nj} - z_{nj}) \mathbf{x}_n.
\end{equation}
需要注意的是，Softmax 回归并不只输出一个硬标签，而是先输出一整组类别概率。如果最高概率只有 0.42 而第二高为 0.39，说明模型并不确定，后续决策时可以保留这种不确定性。

相关图示 的作用是说明 Softmax 并不是只给出一个类别标签，而是先为每个类别计算概率，再选择概率最大的类别作为预测结果。在二维特征空间中，不同类别概率最大的区域共同形成多分类决策边界。

% 【图示占位】源稿图示待按本书风格重绘。

\subsection{支持向量机}

支持向量机（SVM）是一种二分类模型，核心思想是在特征空间中找到一个超平面（Hyperplane），将不同类别的样本分隔开，并使这个超平面到最近样本点的距离（即间隔，Margin）最大化。统计学习理论表明，拥有最大间隔的超平面具有最好的泛化能力。

\paragraph{超平面与间隔}

假设给定数据集 $\mathcal{X}=\{(\mathbf{x}_{n},y_{n})\}_{n=1}^{N}$，其中 $y_n \in \{+1, -1\}$。线性可分 SVM 用一个线性方程来划分数据：
\begin{equation}
    \mathbf{w}^{\textrm{T}}\mathbf{x} + w_0 = 0.
\end{equation}
对于任意数据点 $\mathbf{x}$，其到超平面的距离为 $\frac{|\mathbf{w}^{\textrm{T}}\mathbf{x} + w_0|}{||\mathbf{w}||}$。通过缩放 $\mathbf{w}$ 和 $w_0$，使得距离超平面最近的样本点满足 $|\mathbf{w}^{\textrm{T}}\mathbf{x} + w_0| = 1$，这些点被称为支持向量。几何间隔的总宽度为：
\begin{equation}
    \text{Margin} = \frac{2}{||\mathbf{w}||}.
\end{equation}
因此，最大化间隔等价于最小化 $\frac{1}{2}||\mathbf{w}||^2$。

相关图示 展示了硬间隔 SVM 的理想情形：两类样本完全线性可分，实线是最终分类超平面，两条虚线分别表示 $\mathbf{w}^{\textrm{T}}\mathbf{x}+w_0=\pm 1$。被圈出的样本点就是支持向量，它们直接决定了间隔的位置和宽度；远离边界的样本虽然参与训练，但对最终分隔面的位置影响较小。

% 【图示占位】源稿图示待按本书风格重绘。

\paragraph{优化问题：对偶推导与支持向量}

为了正确分类所有样本并最大化间隔，需要求解以下凸二次规划问题：
\begin{equation}
\begin{aligned}
    \min_{\mathbf{w}, w_0} \quad & \frac{1}{2} ||\mathbf{w}||^2 \\
    \textrm{s.t.} \quad & y_n (\mathbf{w}^{\textrm{T}} \mathbf{x}_n + w_0) \geq 1, \quad \forall n.
\end{aligned}
\end{equation}

\textbf{拉格朗日对偶（Lagrangian Dual）。}引入拉格朗日乘子 $\alpha_n \ge 0$，构造拉格朗日函数：
\begin{equation}
    L(\mathbf{w}, w_0, \mathbf{\alpha}) = \frac{1}{2}||\mathbf{w}||^2 - \sum_{n=1}^{N}\alpha_n [y_n(\mathbf{w}^{\textrm{T}}\mathbf{x}_n + w_0) - 1].
\end{equation}
令 $L$ 对 $\mathbf{w}$ 和 $w_0$ 的偏导为 0，得到 $\mathbf{w} = \sum_{n=1}^{N} \alpha_n y_n \mathbf{x}_n$ 和 $\sum_{n=1}^{N} \alpha_n y_n = 0$。将其代回后，得到对偶问题：
\begin{equation}
\begin{aligned}
    \max_{\alpha} \quad & \sum_{n=1}^{N} \alpha_n - \frac{1}{2}\sum_{i=1}^{N}\sum_{j=1}^{N} \alpha_i \alpha_j y_i y_j \mathbf{x}_i^{\textrm{T}}\mathbf{x}_j \\
    \textrm{s.t.} \quad & \sum_{n=1}^{N} \alpha_n y_n = 0, \quad \alpha_n \ge 0.
\end{aligned}
\end{equation}
对偶问题的一个重要优势是它只涉及样本之间的内积 $\mathbf{x}_i^{\textrm{T}}\mathbf{x}_j$，这为引入核函数奠定了基础。

\textbf{KKT 条件与支持向量。}最优解必须满足互补松弛条件 $\alpha_n [y_n(\mathbf{w}^{\textrm{T}}\mathbf{x}_n + w_0) - 1] = 0$，这意味着只有位于间隔边界上的样本点（支持向量）的 $\alpha_n$ 才可能不为 0，其余样本点对最终分隔面没有影响。

\paragraph{SVM 手算示例}

给定正类样本 $\mathbf{x}_1=(3,3)^\top, \mathbf{x}_2=(4,3)^\top$，负类样本 $\mathbf{x}_3=(1,1)^\top$。需要最小化 $\frac{1}{2}(w_1^2 + w_2^2)$，满足以下约束：
\begin{align*}
    1 - 3w_1 - 3w_2 - w_0 &\le 0 \\
    1 - 4w_1 - 3w_2 - w_0 &\le 0 \\
    1 + w_1 + w_2 + w_0 &\le 0.
\end{align*}
通过求解 KKT 条件方程组，得到最优解 $w_1 = w_2 = 0.5$，$w_0 = -2$，支持向量为 $\mathbf{x}_1$ 和 $\mathbf{x}_3$。最终分隔超平面方程为 $0.5 x_1 + 0.5 x_2 - 2 = 0$。

\paragraph{软间隔与正则化}

现实数据往往充满噪声，并不总是线性可分的。为了允许少量的分类错误，引入松弛变量 $\xi_n \ge 0$，新的优化目标变为：
\begin{equation}
    \min_{\mathbf{w}, w_0} \frac{1}{2}||\mathbf{w}||^2 + C \sum_{n=1}^{N} \xi_n,
\end{equation}
约束条件变为 $y_n (\mathbf{w}^{\textrm{T}} \mathbf{x}_n + w_0) \geq 1 - \xi_n$。其中 $C$ 是正则化参数：$C$ 越大，对误分类的惩罚越大，模型趋向于硬间隔（容易过拟合）；$C$ 越小，容忍更多误分类，间隔越宽（泛化能力可能更强）。实际建模时，$C$ 通常通过验证集或交叉验证选择。

% 【图示占位】源稿图示待按本书风格重绘。

相关图示 中，位于间隔边界之外的样本 $\xi_n=0$；进入两条间隔边界之间的样本 $0<\xi_n<1$；落到分类超平面错误一侧的样本通常有 $\xi_n\geq 1$。

\paragraph{核函数}

对于非线性数据，SVM 通过核技巧（Kernel Trick）将原始低维空间的样本 $\mathbf{x}$ 隐式映射到高维特征空间，使其在高维空间线性可分。核函数 $K(\mathbf{x}_i, \mathbf{x}_j)$ 允许直接在低维空间计算高维内积，常见核函数包括：
\begin{itemize}
    \item 多项式核：$K(\mathbf{x}, \mathbf{x}') = (\mathbf{x}^T \mathbf{x}' + 1)^q$；
    \item 高斯径向基核（RBF）：$K(\mathbf{x}, \mathbf{x}') = \exp \left( - \frac{||\mathbf{x} - \mathbf{x}'||^2}{2 s^2} \right)$；
    \item Sigmoid 核：$K(\mathbf{x}, \mathbf{x}') = \tanh (2 \mathbf{x}^T \mathbf{x}' + 1)$。
\end{itemize}

\paragraph{实现占位}
我们将使用 \texttt{scikit-learn} 中的 \texttt{SVC} 类来演示线性 SVM 和非线性（RBF 核）SVM 的使用。

\begin{quote}
\textbf{【实现占位】} 本段原有代码不直接迁入；后续案例将以 AI 协作过程、关键决策和验证结果为主。
\end{quote}

SVM 在处理中小型数据集、高维数据和具有清晰决策边界的问题时表现优异，核技巧使其能够处理复杂的非线性问题。但其训练时间在非常大的数据集上可能较长，并且对参数选择敏感。

\subsection{决策树}

决策树是一种非参数的监督学习方法，既可以用于分类也可以用于回归。其核心思想采用"分而治之"（divide-and-conquer）的策略，通过一系列递归的划分步骤将特征空间分割为多个区域。决策树由两种节点组成：\textbf{内部决策节点（Internal decision node）}包含一个测试函数，基于特征划分出多个分支；\textbf{叶节点（Terminal leaf node）}定义了特征空间的一个局部区域，落入该区域的实例具有相同的预测输出。

相关图示 可以从特征空间划分的角度理解决策树：每一次节点测试都相当于在特征空间中画出一条划分边界，多个节点递归叠加后形成若干个局部区域。与线性模型用一条整体边界分割空间不同，决策树通过多次轴向切分形成阶梯状区域，解释性强，但树过深时容易把训练数据中的局部噪声切分出来。

% 【图示占位】源稿图示待按本书风格重绘。

\paragraph{分类树的构建与纯度度量}

构建决策树的目标是通过一系列分裂使节点随着树的加深变得更加"纯净"。常用的纯度度量包括三种：分类错误率 $E = 1 - \max_k \hat{p}_{mk}$；交叉熵（信息增益）
\begin{equation}
    D = - \sum_{k=1}^K \hat{p}_{mk} \log \hat{p}_{mk};
\end{equation}
以及基尼指数
\begin{equation}
    G = 1 - \sum_{k=1}^K \hat{p}_{mk}^2,
\end{equation}
其中 $\hat{p}_{mk}$ 是第 $m$ 个区域中属于第 $k$ 类的样本比例。在分裂时，算法尝试所有属性，选择能使纯度增加最大的属性作为测试函数。

对于回归树，构建目标是最小化残差平方和：
\begin{equation}
    \textrm{RSS} = \sum_{j=1}^J \sum_{i \in R_j} (y_i - \hat{y}_{R_j})^2,
\end{equation}
其中 $\hat{y}_{R_j}$ 是第 $j$ 个区域内训练观测值的平均响应。

\paragraph{过拟合与剪枝}

完全生长的树容易过拟合训练数据，导致泛化能力差。解决策略包括两类：\textbf{预剪枝（Pre-pruning）}在树生长过程中提前停止，例如限制最大树深度或节点分裂所需的最小样本数；\textbf{后剪枝（Post-pruning）}先生成一棵非常大的树，然后利用代价复杂度剪枝（Cost Complexity Pruning）将其剪成子树，通过参数 $\alpha$ 权衡树的复杂度与对训练数据的拟合程度。

% 【图示占位】源稿图示待按本书风格重绘。

决策树的主要优点是易于解释、可视化效果好，能处理定性预测变量而无需虚拟变量。主要缺点是预测精度通常低于其他复杂模型，且对数据的微小变化敏感，可能导致树结构发生巨大变化。

\paragraph{实现占位}
我们将使用 \texttt{scikit-learn} 中的 \texttt{DecisionTreeClassifier} 模型，并利用 \texttt{plot\_tree} 进行可视化。

\begin{quote}
\textbf{【实现占位】} 本段原有代码不直接迁入；后续案例将以 AI 协作过程、关键决策和验证结果为主。
\end{quote}

\subsection{随机森林}

随机森林是一种集成学习方法，通过构建大量决策树并将它们的预测结果进行平均（回归）或投票（分类）来提高模型的准确性和鲁棒性。随机森林的构建包含两个关键的随机性来源：\textbf{自助采样法（Bagging / Bootstrap Aggregating）}从原始训练数据集中有放回地随机抽取样本，生成多个不同的子数据集，每个子数据集用于训练一棵独立的决策树；\textbf{特征随机性}在每个节点进行分裂时随机选择特征的一个子集，增加树之间的多样性。

通过这些随机性，随机森林中的每棵树都具有一定的独立性，从而减少了单棵决策树的过拟合风险，提高了模型的泛化能力。随机森林还能提供特征重要性评估，帮助我们理解哪些特征对预测最关键。缺点是模型解释性不如单棵决策树直观。

\paragraph{实现占位}
我们将使用 \texttt{scikit-learn} 中的 \texttt{RandomForestClassifier} 模型。

\begin{quote}
\textbf{【实现占位】} 本段原有代码不直接迁入；后续案例将以 AI 协作过程、关键决策和验证结果为主。
\end{quote}

\section{无监督学习模型}

无监督学习处理的是没有标签的数据，即只有特征 $\mathbf{x}$ 而没有标签 $y$。其目标是发现数据中固有的结构、规律和模式。常见的应用包括聚类（将相似的样本分组）和降维（发现能够代表原始数据大部分信息的潜在结构）。与监督学习相比，无监督学习的评估往往更困难，因为没有"标准答案"可以直接对比。

\subsection{K-均值聚类}

K-均值聚类旨在将数据集 $\{\mathbf{x}_1, \ldots, \mathbf{x}_N\}$ 划分为 $K$ 个簇，目标是找到 $K$ 个簇中心 $\boldsymbol{\mu}_k$ 并分配每个数据点到最近的簇中心，使得簇内平方和（Inertia）最小化：
\begin{equation}
    J = \sum_{n=1}^N \sum_{k=1}^K r_{nk} ||\mathbf{x}_n - \boldsymbol{\mu}_k||^2,
\end{equation}
其中 $r_{nk} \in \{0, 1\}$ 是二值指示变量，$r_{nk}=1$ 表示数据点 $\mathbf{x}_n$ 被分配给簇 $k$。

\paragraph{算法流程}

K-Means 算法是一个迭代优化过程：初始化 $K$ 个簇中心 $\boldsymbol{\mu}_k$；在分配阶段（E-step 思想）固定 $\boldsymbol{\mu}_k$，将每个数据点分配给最近的簇中心；在更新阶段（M-step 思想）固定 $r_{nk}$，重新计算各簇中心；重复直到收敛。

在更新阶段，对 $\boldsymbol{\mu}_k$ 求偏导并令其为 0，得到：
\begin{equation}
    \boldsymbol{\mu}_k = \frac{\sum_{n=1}^N r_{nk} \mathbf{x}_n}{\sum_{n=1}^N r_{nk}}.
\end{equation}
这证明了在平方误差准则下，簇内所有点的算术平均值确实是使误差最小的最优质心。K-Means 只能保证收敛到局部较优解，因此实际使用中常会多次初始化并选择 Inertia 最小的一组结果。

相关图示 展示了 K-Means 的两个交替步骤：先根据当前质心把样本分配到最近簇，再用每个簇内样本的平均值更新质心，不断降低簇内平方误差直到收敛。

% 【图示占位】源稿图示待按本书风格重绘。

\paragraph{K 值的选择：手肘法}

在无监督学习中，没有直接的误差指标来评估聚类效果。手肘法通过绘制 Inertia 随 $K$ 变化的曲线来确定最佳簇数：最佳点通常位于曲线的"肘部"，即 Inertia 下降速度显著变缓的点。阅读相关图示时，重点不是寻找最低点（因为 $K$ 越大 Inertia 通常越小），而是观察下降幅度从"明显变好"转为"收益有限"的位置。该位置对应的 $K$ 往往能在模型复杂度与聚类解释性之间取得较好平衡。

% 【图示占位】源稿图示待按本书风格重绘。

\paragraph{实现占位}
\begin{quote}
\textbf{【实现占位】} 本段原有代码不直接迁入；后续案例将以 AI 协作过程、关键决策和验证结果为主。
\end{quote}

K-均值聚类算法简单、高效，适用于大数据集。但它需要预先指定簇的数量 $K$，对初始簇中心的选择敏感，并且对噪声和异常值比较敏感。此外，它假设簇是球形的且大小相似，对于非球形或密度差异大的簇效果不佳。

\subsection{高斯混合模型}

\paragraph{概率论回顾：贝叶斯定理}

在介绍高斯混合模型（GMM）之前，先回顾贝叶斯定理，它是概率聚类的基础。贝叶斯法则描述了先验概率（Prior）和后验概率（Posterior）之间的关系：
\begin{equation}
    P(A|B) = \frac{P(B|A)P(A)}{P(B)}.
\end{equation}

\textbf{应用实例：医疗诊断。}假设人群中某种疾病的患病率（先验概率）为 $P(A) = 0.001$，某检测方法的敏感性（真阳性率）为 95\%，特异性（真阴性率）为 98\%。若检测结果呈阳性，真正患病的概率 $P(A|B)$ 是多少？根据贝叶斯定理：
\begin{align}
    P(A|B) = \frac{0.95 \times 0.001}{0.95 \times 0.001 + 0.02 \times 0.999} \approx 0.045.
\end{align}
即使检测呈阳性，由于先验患病率极低，患者真正患病的概率仅为 4.5\%。这说明了先验概率（即 GMM 中的混合系数 $\pi_k$）对最终归类的巨大影响。

K-均值聚类是一种硬聚类（Hard Clustering），每个点只能属于一个簇。而高斯混合模型是 K-Means 的概率扩展，属于软聚类（Soft Clustering），给出了数据点属于某个簇的概率。GMM 假设数据是由 $K$ 个高斯分布混合生成的，概率密度函数为：
\begin{equation}
    p(\mathbf{x}) = \sum_{k=1}^K \pi_k \mathcal{N}(\mathbf{x} | \boldsymbol{\mu}_k, \boldsymbol{\Sigma}_k),
\end{equation}
其中 $\pi_k$ 是混合系数（先验概率），满足 $\sum \pi_k = 1$；$\mathcal{N}(\mathbf{x} | \boldsymbol{\mu}_k, \boldsymbol{\Sigma}_k)$ 是第 $k$ 个高斯分量的密度函数。

% 【图示占位】源稿图示待按本书风格重绘。

\paragraph{EM 算法 (Expectation-Maximization)}

GMM 的参数 $\{\pi_k, \boldsymbol{\mu}_k, \boldsymbol{\Sigma}_k\}$ 无法直接通过最大似然估计求解，通常使用 EM 算法迭代求解：
\begin{enumerate}
    \item \textbf{E 步（Expectation）}：计算后验概率（Responsibility）$\gamma(z_{nk})$，即给定当前参数下，数据点 $n$ 由第 $k$ 个高斯分量生成的概率：
    \begin{equation}
        \gamma(z_{nk}) = \frac{\pi_k \mathcal{N}(\mathbf{x}_n | \boldsymbol{\mu}_k, \boldsymbol{\Sigma}_k)}{\sum_{j=1}^K \pi_j \mathcal{N}(\mathbf{x}_n | \boldsymbol{\mu}_j, \boldsymbol{\Sigma}_j)}.
    \end{equation}
    \item \textbf{M 步（Maximization）}：利用 E 步得到的 $\gamma(z_{nk})$ 更新参数（令 $N_k = \sum_{n=1}^N \gamma(z_{nk})$）：
    \begin{align*}
        \boldsymbol{\mu}_k^{new} &= \frac{1}{N_k} \sum_{n=1}^N \gamma(z_{nk}) \mathbf{x}_n, \qquad
        \boldsymbol{\Sigma}_k^{new} = \frac{1}{N_k} \sum_{n=1}^N \gamma(z_{nk}) (\mathbf{x}_n - \boldsymbol{\mu}_k^{new})(\mathbf{x}_n - \boldsymbol{\mu}_k^{new})^T, \qquad
        \pi_k^{new} = \frac{N_k}{N}.
    \end{align*}
\end{enumerate}

\paragraph{模型选择：如何确定 GMM 的簇数}

与 K-Means 使用 Inertia 不同，GMM 使用基于似然函数的统计量来选择最佳的簇数 $K$。常用指标包括赤池信息量准则（AIC）和贝叶斯信息量准则（BIC）：
\begin{align}
    \text{AIC} = 2p - 2\ln(L), \qquad \text{BIC} = p\ln(N) - 2\ln(L),
\end{align}
其中 $N$ 是样本数量，$p$ 是参数数量，$L$ 是似然函数值。较小的 AIC 或 BIC 值表示模型更好；BIC 对模型复杂度的惩罚比 AIC 更重，倾向于选择更简单的模型。

\paragraph{工程应用案例：基于 GMM 的交通速度分布建模}

GMM 在土木工程中常用于处理异质性数据。例如，Park 等人（2010）使用贝叶斯混合模型来捕捉车辆速度数据的异质性。在美国得克萨斯州 IH-35 高速公路上收集的 24 小时交通数据中，交通流在不同时间段表现出不同模式（如自由流、拥堵流），单一的高斯分布无法拟合整体速度分布。通过 $K=4$ 的混合模型，可以清晰地识别出不同交通状态下的速度分布模式（如夜间高速模式、早高峰低速模式等），这比单一分布提供了更丰富的信息，研究者通过计算 AIC/BIC 确定了 $K=4$ 是最优选择。

% 【图示占位】源稿图示待按本书风格重绘。

相关图示 中，每一条红色虚线可以理解为一种典型交通状态下的速度分布，蓝色实线则是这些状态按权重叠加后的总体分布。混合模型比用单个高斯分布更能解释交通系统在自由流、过渡流和拥堵流之间切换的异质性。

\paragraph{实现占位}
\begin{quote}
\textbf{【实现占位】} 本段原有代码不直接迁入；后续案例将以 AI 协作过程、关键决策和验证结果为主。
\end{quote}

GMM 比 K-均值更灵活，能够更好地处理非球形簇和不同密度的簇，并提供每个数据点属于每个簇的概率。缺点是计算成本更高，对初始值也相对敏感（EM 算法容易陷入局部最优），且需要选择合适的协方差类型和分量数量。

\subsection{降维：特征选择与主成分分析}

在高维数据中，特征数量往往影响学习算法的时间和空间复杂度，甚至导致过拟合或统计性质变差（维度诅咒）。降维主要有两种方式：\textbf{特征选择（Feature Selection）}直接从原始特征中选择 $k < d$ 个重要特征，丢弃其余特征，保留特征的物理意义；\textbf{特征提取（Feature Extraction）}如 PCA，将原始 $d$ 维空间投影到新的 $k$ 维空间，新特征是原始特征的线性组合。

\paragraph{特征选择}

特征选择的目标是找到一个最佳特征子集，使特定任务的准确率最高。由于遍历所有 $2^d - 1$ 个非空子集通常是不可行的（NP-hard），常使用启发式算法。

\textbf{前向搜索（Forward Search）}从空集 $F = \emptyset$ 开始：每轮迭代对每个可用特征 $x_i$ 训练模型并计算验证集误差 $E(F \cup x_i)$，选择使误差最小的特征 $x_j$ 加入 $F$，直到没有显著改进为止。

\textbf{后向搜索（Backward Search）}从包含所有特征的集合开始：每轮迭代尝试移除对模型贡献最小的特征，直到移除会导致模型显著退化为止。

\paragraph{主成分分析（PCA）}

主成分分析是一种常用的特征提取技术。它通过正交变换，将可能相关的变量转换为线性不相关的主成分，第一个主成分拥有最大的方差（包含最多的信息），主成分是协方差矩阵 $\boldsymbol{\Sigma}$ 的特征向量。

\textbf{PCA 的数学推导。}PCA 的目标是找到投影方向 $\mathbf{w}_1$，使投影后数据方差 $\text{Var}(z_1) = \mathbf{w}_1^T \boldsymbol{\Sigma} \mathbf{w}_1$ 最大，且满足 $||\mathbf{w}_1|| = 1$。构造拉格朗日函数 $L = \mathbf{w}_1^T \boldsymbol{\Sigma} \mathbf{w}_1 - \alpha (\mathbf{w}_1^T \mathbf{w}_1 - 1)$，对 $\mathbf{w}_1$ 求导并令其为 0，得到特征值方程：
\begin{equation}
    \boldsymbol{\Sigma} \mathbf{w}_1 = \alpha \mathbf{w}_1.
\end{equation}
这意味着 $\mathbf{w}_1$ 必须是协方差矩阵 $\boldsymbol{\Sigma}$ 的特征向量，$\alpha$ 是对应的特征值。为了最大化方差，需要选择最大特征值对应的特征向量作为第一主成分。

\textbf{方差解释比例（PoV）。}为了量化降维保留了多少信息，计算前 $k$ 个主成分解释的方差比例：
\begin{equation}
    PoV = \frac{\sum_{i=1}^k \lambda_i}{\sum_{j=1}^d \lambda_j},
\end{equation}
其中 $\lambda_i$ 是协方差矩阵的特征值（按降序排列）。通常选择截断点 $k$，使得 $PoV > 0.9$（即保留 90\% 的原始信息）。碎石图（Scree Graph）用于直观确定保留多少个主成分：若前几个主成分已经解释了大部分方差，就可以用较低维表示替代原始高维特征。

相关图示 的关键是"换坐标系"：PCA 不是随便丢掉几个原始变量，而是先寻找数据变化最明显的方向，把它定义为第一主成分，再在正交方向中寻找第二主成分。样本投影到主成分方向后，原本由多个相关特征共同表达的信息，可以被更少的新坐标概括。

% 【图示占位】源稿图示待按本书风格重绘。

\paragraph{实现占位}
\begin{quote}
\textbf{【实现占位】} 本段原有代码不直接迁入；后续案例将以 AI 协作过程、关键决策和验证结果为主。
\end{quote}

特征选择通过保留物理意义明确的特征来简化模型，而 PCA 通过正交变换提取主要信息。PCA 是一种强大的线性降维技术，但它是无监督的，可能会丢失对分类重要的信息，且无法处理复杂的非线性流形结构。

\section{案例占位}

\begin{quote}
\textbf{【案例占位】}

后续将在此加入一个围绕。案例不以逐行教学代码为主，而将展示问题定义、向 AI 提供的上下文与约束、模型选择依据、关键中间产物、验证过程以及人需要承担的判断。

\textbf{【待确定内容】} 数据来源、任务背景、AI 协作流程、模型比较、错误诊断与现实解释。
\end{quote}
